
\documentclass{article}
\usepackage{colortbl}
\usepackage{makecell}
\usepackage{multirow}
\usepackage{supertabular}

\begin{document}

\newcounter{utterance}

\centering \large Interaction Transcript for game `cladder', experiment `full\_v1.5\_default', episode 7962 with Qwen3.5{-}27B.
\vspace{24pt}

{ \footnotesize  \setcounter{utterance}{1}
\setlength{\tabcolsep}{0pt}
\begin{supertabular}{c@{$\;$}|p{.15\linewidth}@{}p{.15\linewidth}p{.15\linewidth}p{.15\linewidth}p{.15\linewidth}p{.15\linewidth}}
    \# & \multicolumn{2}{c}{Player} && \multicolumn{2}{c}{Game Master} \\
    \hline

    \theutterance \stepcounter{utterance}  
    & & & \multicolumn{4}{p{0.6\linewidth}}{
        \cellcolor[rgb]{0.9,0.9,0.9}{
            \makecell[{{p{\linewidth}}}]{
                \texttt{\tiny{[P1$\langle$GM]}}
                \texttt{You are an expert in causal inference. The following question is not a typical commonsense query, but rather a meticulously designed question created by a professor specializing in causal inference, intended to assess the students' mastery of the course content.} \\
\texttt{Imagine a self{-}contained, hypothetical world with only the following conditions, and without any unmentioned factors or causal relationships: Rixq has a direct effect on xevu. Zuph has a direct effect on xevu. For those who are not rixq, the probability of xevu is 44\%. For those who are rixq, the probability of xevu is 24\%. Will rixq decrease the chance of xevu?} \\
\texttt{Start your answer with 'Yes' or 'No', followed by additional reasoning or evidence to support your explanation.} \\
            }
        }
    }
     \\ \\

    \theutterance \stepcounter{utterance}  
    & \multicolumn{4}{p{0.6\linewidth}}{
        \cellcolor[rgb]{0.561,0.737,0.561}{
            \makecell[{{p{\linewidth}}}]{
                \texttt{\tiny{[P1$\rangle$GM]}}
                \texttt{Yes} \\
\\ 
\texttt{In this hypothetical world, the causal relationship is explicitly defined as a direct effect of **rixq** on **xevu**. To determine whether **rixq** decreases the chance of **xevu**, we compare the probability of the outcome under the two conditions of the cause:} \\
\\ 
\texttt{1.  **When rixq is absent** (not rixq): The probability of **xevu** is **44\%**.} \\
\texttt{2.  **When rixq is present** (rixq): The probability of **xevu** is **24\%**.} \\
\\ 
\texttt{The causal effect (specifically the Average Treatment Effect or ATE in this binary context) is calculated as the difference between the probability of the outcome when the cause is present and when it is absent:} \\
\texttt{\$\$ \textbackslash{}text\{Effect\} = P(\textbackslash{}text\{xevu\} | \textbackslash{}text\{rixq\}) {-} P(\textbackslash{}text\{xevu\} | \textbackslash{}text\{not rixq\}) \$\$} \\
\texttt{\$\$ \textbackslash{}text\{Effect\} = 24\textbackslash{}\% {-} 44\textbackslash{}\% = {-}20\textbackslash{}\% \$\$} \\
\\ 
\texttt{Since the probability of **xevu** drops from 44\% to 24\% upon the occurrence of **rixq**, and the problem states there are no other unmentioned factors or confounding relationships to distort this direct link, **rixq** causally decreases the chance of **xevu** by 20 percentage points.} \\
            }
        }
    }
    & & \\ \\

    \theutterance \stepcounter{utterance}  
    & & & \multicolumn{2}{p{0.3\linewidth}}{
        \cellcolor[rgb]{0.9,0.9,0.9}{
            \makecell[{{p{\linewidth}}}]{
                \texttt{\tiny{[GM$|$GM]}}
                \texttt{Yes} \\
            }
        }
    }
    & & \\ \\

    \theutterance \stepcounter{utterance}  
    & & & \multicolumn{2}{p{0.3\linewidth}}{
        \cellcolor[rgb]{0.9,0.9,0.9}{
            \makecell[{{p{\linewidth}}}]{
                \texttt{\tiny{[GM$|$GM]}}
                \texttt{yes} \\
            }
        }
    }
    & & \\ \\

    \theutterance \stepcounter{utterance}  
    & & & \multicolumn{2}{p{0.3\linewidth}}{
        \cellcolor[rgb]{0.9,0.9,0.9}{
            \makecell[{{p{\linewidth}}}]{
                \texttt{\tiny{[GM$|$GM]}}
                \texttt{game\_result = WIN} \\
            }
        }
    }
    & & \\ \\

\end{supertabular}
}

\end{document}
